\chapter{Ensayos e inspección de materiales}\label{ch:11}
\objetivos{Elegir una técnica a partir del defecto y el material; distinguir indicación, defecto y criterio de aceptación; interpretar el alcance real de una inspección negativa.}
\section{Medir propiedades e inspeccionar una pieza}
Los ensayos mecánicos caracterizan respuestas bajo condiciones de carga. La inspección busca indicios de discontinuidades o degradación en un objeto concreto. Un ensayo de tracción no sustituye una inspección de grietas y una inspección sin indicaciones no mide por sí sola el límite elástico.

La tracción aporta curva y magnitudes definidas; la dureza mide resistencia a penetración bajo método específico; el impacto caracteriza energía absorbida en una geometría y velocidad dadas. Fatiga y fluencia necesitan controlar carga, temperatura, tiempo, entorno y dispersión. Cambiar probeta o condición puede impedir comparar números directamente.
\section{Indicación, evaluación y aceptación}
Una indicación es una respuesta del método que requiere interpretación. Puede ser relevante, no relevante o falsa. Un defecto rechazable es una condición que incumple un criterio aplicable; no todo contraste observado tiene ese significado. La aceptación depende de documentación y ubicación, además del tamaño.

Se documentan método, procedimiento, equipo, calibración, sensibilidad, superficie, accesibilidad, orientación, operador y fecha. La resolución nominal no es probabilidad de detección. Una probabilidad de detección se determina para un sistema de inspección y una población representativa de defectos; no se deduce de un único patrón que el equipo «ve».
\section{Técnicas de superficie}
La inspección visual depende de iluminación, línea de visión, aumento, preparación y experiencia. La boroscopia permite acceso a zonas internas, pero tiene restricciones de perspectiva, cobertura y medida. El aspecto de una grieta no siempre determina su profundidad.

Los líquidos penetrantes resaltan discontinuidades abiertas a una superficie no porosa. Preparación, limpieza, tiempo y compatibilidad son importantes. No detectan una grieta cerrada completamente subsuperficial. Las partículas magnéticas requieren material ferromagnético y una magnetización adecuada para que discontinuidades superficiales o próximas generen indicaciones. La orientación del campo respecto al defecto importa; no se aplica de la misma forma a aluminio o a todos los inoxidables.
\section{Técnicas volumétricas y electromagnéticas}
Los ultrasonidos utilizan propagación y reflexión de ondas; permiten estudiar discontinuidades internas, espesores y ciertas uniones. El acoplamiento, la orientación, la atenuación, la estructura de grano, la geometría y las zonas muertas limitan capacidad. Un defecto desfavorablemente orientado puede ser difícil de detectar con una configuración concreta.

La radiografía responde a cambios en atenuación a lo largo del haz. Puede ser útil para porosidad, inclusiones o geometrías internas, pero una grieta plana de abertura pequeña y orientación desfavorable puede presentar poco contraste. El acceso y la seguridad radiológica condicionan la técnica.

Las corrientes inducidas requieren conductividad eléctrica. Se relacionan con cambios en impedancia por discontinuidades, conductividad, permeabilidad, separación de sonda y geometría. La profundidad de penetración varía con frecuencia y propiedades; no se trata de una técnica universal para grandes profundidades.
\begin{equation}\delta=\frac1{\sqrt{\pi f\mu\sigma_e}}\quad\text{(conductor ideal homogéneo).}\label{eq:skin}\end{equation}
Aquí $\sigma_e$ es conductividad eléctrica, no tensión mecánica. La expresión de profundidad pelicular describe atenuación ideal y no equivale al tamaño mínimo detectable por una sonda real.
\section{Termografía y emisión acústica}
La termografía observa respuesta térmica superficial, en ocasiones tras una excitación controlada. Profundidad, emisividad, tiempo y contraste condicionan detección de despegues o delaminaciones. Una imagen de temperatura no identifica automáticamente la causa del contraste.

La emisión acústica detecta señales transitorias asociadas a actividad de fuentes bajo condiciones de carga. Puede ayudar a localizar actividad, pero requiere discriminar ruido, propagación y mecanismos. Una discontinuidad inactiva puede no emitir; el silencio del sistema no significa ausencia de todos los defectos.
\begin{table}[htbp]\centering\small\caption{Selección inicial de END: capacidades y limitaciones esenciales.}\label{tab:end}
\begin{tabular}{p{.18\linewidth}p{.32\linewidth}p{.38\linewidth}}\toprule
Técnica & Evidencia habitual & Restricción que no debe omitirse\\\midrule
Visual/boroscopia & Alteración visible y accesible & Cobertura, iluminación y medida de profundidad\\[3pt]
Penetrantes & Apertura a superficie & Superficie no porosa y preparación\\[3pt]
Magnéticas & Discontinuidades en ferromagnéticos & Material y orientación de magnetización\\[3pt]
Ultrasonidos & Reflexiones y tiempos de vuelo & Acoplamiento, orientación y atenuación\\[3pt]
Radiografía & Diferencias de atenuación & Superposición y orientación de grietas\\[3pt]
Inducidas & Alteración electromagnética & Conductividad, profundidad y separación\\[3pt]
Termografía & Contraste térmico & Emisividad, profundidad y excitación\\\bottomrule
\end{tabular}\end{table}
\caso{M11. Dos defectos, dos preguntas}{Se busca una grieta abierta en una pieza de aluminio y una delaminación interna en un laminado. Las partículas magnéticas no son una respuesta universal para la primera; se estudian penetrantes o corrientes inducidas según acceso y datos. Para la segunda pueden valorarse ultrasonidos o termografía con un procedimiento adecuado. La selección final exige material, geometría, tamaño y orientación objetivo, más criterios de aceptación. El caso no establece un procedimiento aprobado.}
\section{Cómo cerrar una inspección}
La cadena de registro debe permitir reconstruir qué zona se examinó y con qué sensibilidad. Un informe «sin defectos» sin cobertura ni procedimiento ofrece poca información. La evaluación de repetibilidad necesita comparar configuraciones y condiciones, no solo resultados escritos.

Para estimar intervalos se combina la probabilidad de detección con un modelo de crecimiento y márgenes. La resolución y el tamaño crítico no son por sí solos una política de inspección. La discusión de fractura y fatiga de los capítulos~\ref{ch:03} y~\ref{ch:04} proporciona el contexto físico de esa decisión.
\section{Elección de ensayos no destructivos en preguntas de teoría}
La técnica se decide con la información del material, profundidad, orientación y tipo de discontinuidad, geometría, acceso, capacidad de detección y procedimiento aplicable. No existe un método universal para medir cualquier grieta de cinco milímetros. Penetrantes: discontinuidades abiertas a superficie no porosa; partículas magnéticas: materiales ferromagnéticos; corrientes inducidas: materiales conductores; ultrasonidos: reflexión y propagación dependientes del recorrido; radiografía: contraste por atenuación y orientación del defecto. Una grieta planar fina desfavorablemente orientada puede ser poco visible en radiografía.

La termografía infrarroja interpreta diferencias de radiación emitida relacionadas con temperatura y emisividad; no mide directamente el tamaño de una discontinuidad. Para pulso--eco ideal en incidencia normal, $d=vt/2$. Una amplitud de eco mayor no prueba automáticamente un defecto mayor: influyen orientación, calibración, geometría, acoplamiento y atenuación.


\section{Detectabilidad, resolución y probabilidad}
La resolución espacial distingue detalles; la detectabilidad depende además de contraste, ruido, orientación, procedimiento y decisión del inspector o sistema. Una indicación menor que una dimensión nominal del equipo puede detectarse en ciertas condiciones, y un defecto mayor puede pasar inadvertido en otras. No se transforma la resolución del equipo en un tamaño garantizado de grieta detectable.

La probabilidad de detección, POD, se estima con una población y un procedimiento de ensayo. Debe acompañarse de incertidumbre y del criterio estadístico usado. También existen falsas llamadas y errores de dimensionado. Un resultado negativo no prueba tamaño cero; se incorpora a un argumento de cobertura y capacidad compatible con el programa de inspección. Un valor $a_{90/95}$, cuando se usa, expresa una convención de probabilidad y confianza, no una certeza individual universal.
\section{Comparación por mecanismo físico}
Penetrantes requiere abertura a una superficie accesible y compatible; una capa que selle la abertura puede ocultar el defecto. Partículas magnéticas necesita material ferromagnético y orientación de campo pertinente. Corrientes inducidas explora perturbaciones electromagnéticas en conductores y está condicionada por frecuencia, geometría, propiedades y lift-off. La profundidad pelicular de P28 describe atenuación ideal, no una garantía de detección.

Ultrasonidos necesita acoplamiento y un recorrido interpretable; interfaces, grano y orientación afectan señal. Radiografía representa atenuación integrada y responde de modo distinto a discontinuidades volumétricas y planos mal orientados; una proyección puede superponer geometrías. Termografía observa respuesta térmica y depende de excitación, profundidad y propiedades. Emisión acústica escucha eventos durante una solicitación, sin proporcionar automáticamente una imagen completa de defectos inactivos.
\ilustra{ultrasonidos}{Modelo pulso--eco: el tiempo de vuelo incluye ida y vuelta. La posición de los ecos es ilustrativa y no constituye una calibración de equipo.}
\Needspace{.48\textheight}
\section{Ensayar una propiedad y decidir aceptación}
Tracción proporciona una curva dependiente de probeta, velocidad, temperatura y medida; dureza es una respuesta local bajo un método concreto. Las correlaciones entre dureza y resistencia son empíricas y no universales. Impacto caracteriza una respuesta dinámica de una geometría y condición de ensayo; su energía no es $K_{IC}$. Fatiga y fluencia añaden historia temporal y dispersión.

Un ensayo o END produce información; la aceptación compara esa información con criterios aplicables y considera cobertura, incertidumbre y configuración. Conservar método, equipo, revisión, zona, resultados y personal permite interpretar el informe años después. Una casilla «OK» sin contexto puede ser insuficiente para reconstruir qué se examinó.
\repaso{¿Por qué dos END que exploran la misma pieza pueden dar resultados diferentes sin que uno sea necesariamente incorrecto?}{Pueden ser sensibles a mecanismos, orientaciones y profundidades distintas, con diferentes coberturas y umbrales. Debe compararse la capacidad física de cada técnica y el procedimiento aplicado antes de atribuir una contradicción.}

\resumen{La técnica se elige desde el defecto, material y acceso. Una indicación no es automáticamente un defecto rechazable. La ausencia de indicaciones solo tiene significado dentro de la cobertura y capacidad del sistema. La calibración no demuestra por sí sola validez para cualquier geometría.}
\fuentes{FAA--General, cap. 10, \textit{Inspection Concepts and Techniques}, y cap. 7 \cite{FAA-G}; materiales docentes de control y detección de fallos.}
